\chapter{Conclusion}
\label{chap-conclusion}

The first version of this project answered the question asked in 2023: do we know how to
build a graphical interface with GTK? It answered it well --- the game was complete --- but
it left aside the other half of the course, the one that talks about organisation: one
concept per module, use independent of implementation, a function written and therefore
tested.

The second version answers that other half. What it demonstrates comes down to three points.

\textbf{Separating costs little and pays a lot.} Extracting the rules into a library with no
dependency made possible, in one stroke, the automatic tests, the second interface, the
measurement of the strategies and reproducibility. None of these four things was conceivable
as long as a cell of the game contained the address of a button.

\textbf{Measuring corrects intuition.} We thought, in 2023, that the strength of the computer
lay in the way it pursued a hit ship. The measurements show that the essential lies
elsewhere: respecting what the no-contact rule proves saves thirty shots, the pursuit fewer
than eight. The difficulty levels were redefined according to these figures.

\textbf{Tools find what use hides.} The out-of-bounds read of the hard level never crashed
the game in front of us. It was there all the same, and a standard tool would have reported
it at the very first run. Verifying automatically, at every change, is not an extra comfort:
it is what makes it possible to modify a program without fear of breaking it.

The game itself has not changed. One still commands the navy of Sup Galilée against the
pirates, on a grid of ten cells by ten. That was the goal: for the player to see no
difference, and for the next developer to see nothing else.
